\section{What Is Proved, What Is Computed, What Remains Open}
\label{sec:proved-supported-deferred}

\Cref{tab:claim-support-matrix} summarizes the status of each claim. This section comments on the three groups.

\begin{table}[t]
    \caption{Status of the main claims. ``Proved'' means proved in the paper and, where indicated, in Lean. ``Cited'' means a classical result used as stated in the literature. ``Computed'' means verified on the specific systems examined. ``Open'' means not established here.}
    \label{tab:claim-support-matrix}
    \centering
    \smallskip
    \small
\setlength{\tabcolsep}{5pt}
\renewcommand{\arraystretch}{1.15}
\begin{tabularx}{\linewidth}{@{}Y L{0.30\linewidth} l@{}}
\toprule
Claim & Evidence & Status \\
\midrule
Elementary flatness $\iff$ local confluence (\Cref{thm:elementary-flatness-iff-local-confluence}) & paper proof; Lean & proved \\
Confluence and normalization give a unique normal form (\Cref{prop:confluence-normalization-unique-normal-form}) & paper proof; Lean & proved \\
Newman's lemma; under termination, confluence $\iff$ elementary flatness (\Cref{thm:newman}, \Cref{cor:terminating-confluence-iff-flat}) & paper proof; Lean & proved \\
Under normalization, confluence $\iff$ unique reachable normal forms (\Cref{thm:normalizing-confluence-iff-unf}) & paper proof; Lean & proved \\
Cycle-based graph flatness does not detect confluence (\Cref{prop:bare-graph-not-diagnostic}) & explicit counterexamples & proved \\
Critical-pair lemma and its holonomy form (\Cref{thm:critical-pair-lemma}, \Cref{cor:cp-flatness}) & literature \cite{KnuthBendix1970,Huet1980} & cited \\
Theorems agree with all terminating systems on $1$--$4$ states; each hypothesis needed (\Cref{tab:exhaustive-finite-ars-summary,tab:boundary-counterexamples}) & exhaustive enumeration & computed \\
Reachable critical-pair instances match the observed obstructions in the curated string and term examples (\Cref{prop:reachable-critical-pairs-capture-string-obstructions,prop:trs-critical-pairs-as-curvature-generators}) & complete exploration of curated examples & computed \\
A single added rule fills the open branching in two curated pairs (\Cref{prop:completion-as-holonomy-elimination}) & complete exploration of curated examples & computed \\
Mechanized critical-pair lemma for left-linear term rewriting in 2-complex form & --- & open \\
Group-valued or connection-based holonomy theory of confluence & --- & open \\
\bottomrule
\end{tabularx}

\end{table}

\subsection{Proved}
\label{subsec:what-is-proved}

For abstract rewriting systems the picture is complete. Elementary flatness is local confluence (\Cref{thm:elementary-flatness-iff-local-confluence}). Under termination, Newman's lemma makes it equivalent to confluence (\Cref{cor:terminating-confluence-iff-flat}); under the weaker hypothesis of normalization, confluence is equivalent to uniqueness of reachable normal forms (\Cref{thm:normalizing-confluence-iff-unf}); and a terminating, elementarily flat system has exactly one normal form reachable from each state (\Cref{cor:terminating-flat-unique-nf}). These are classical facts of rewriting theory restated in the vocabulary of holonomy, and all of them are checked in Lean for arbitrary, possibly infinite, systems. The failure of cycle-based graph flatness (\Cref{prop:bare-graph-not-diagnostic}) is proved by explicit counterexamples.

\subsection{Cited and computed}
\label{subsec:what-is-supported}

For string and term rewriting, the link between critical pairs and local confluence is the classical critical-pair lemma, which we cite (\Cref{thm:critical-pair-lemma}). Together with the results above it gives \Cref{cor:cp-flatness}: a system is elementarily flat if and only if its critical pairs are joinable, and, when terminating, confluent under the same condition.

The computations add evidence of a different kind. On finite systems with one to four states they agree with the theorems on every case and produce smallest counterexamples for each dropped hypothesis. On the curated string and term examples they show that the start-relative analysis, based on concrete critical-pair instances in reachable terms, finds exactly the observed obstructions, including an overlap below the root, and that a single added rule can turn an open branching into a filled one. These are statements about the examples examined, and we do not extrapolate them.

\subsection{Open}
\label{subsec:what-remains}

Three directions remain open.

\begin{enumerate}
    \item \emph{A mechanized critical-pair lemma.} The Lean development covers abstract rewriting only. A formal proof of \Cref{thm:critical-pair-lemma} and its start-relative form for left-linear term rewriting, phrased in terms of reduction 2-complexes, requires the full case analysis of \Cref{subsec:three-kinds-of-peaks} together with stability of rewriting under contexts and substitutions. It is not attempted here.
    \item \emph{A genuinely geometric formulation.} Elementary holonomy is a Boolean status. A theory in which holonomy takes values in a group, or in which confluence becomes the flatness of an independently defined connection on a reduction complex, would need definitions that this paper does not supply. The homotopical approach to rewriting \cite{Squier1987,GuiraudMalbos2012,GuiraudMalbos2018,Mimram2023} is the natural setting for such a development.
    \item \emph{Completion in general.} \Cref{sec:completion-holonomy} reads two single completion steps as holonomy elimination. Whether the full completion procedure, with its orderings, new critical pairs, and possible failure, admits a useful description in these terms is open.
\end{enumerate}
